%! root=../main.tex
\section{Conclusion}

We present \pname{}, the first formally verified binary CFI rewriter implementation for a mainstream ISA\@.
Verifying binary rewriters is challenging as the verifier must reason about non-concrete code.
\pname{} addresses this challenge through a dependently-typed, coinductive proof strategy based on Rocq and \Picinae.
We present formal verification-friendly design principles to guide future binary rewriter designs, and present case studies of test-evading bugs that were caught in thanks to program-proof co-development.

\comment{
    Spot checked 4 works from CCS'25 Formal Methods+PL track
    
    * https://dl.acm.org/doi/epdf/10.1145/3719027.3765213
        Discussion only, no separate conclusion
        Starts off with the typical conclusion: In this paper we did X written in Y. The key features are X and Y rather than Z. In this way, were were able to prove P. Our impl. and proofs submitted as artifact.
        The results are summarized in Table. We used three tools, X for A, Y for B, Z for C. Using a single tool would be better because X, but we decide against it for pragmatic reasons.
    
        And then it has two bolded sections: Comparison with Other Approaches (in lieu of related works) and Future Directions.
    
    * https://dl.acm.org/doi/epdf/10.1145/3719027.3765104
        Conclusion is 1.2 paragraphs: We did X. A key success factor was Y, which allowed A, B, C. ``We discovered both design and implementation errors, critically affecting the router’s security, that evaded all prior reviews and testing and have all been subsequently fixed''.
    
        Second par: Future work includes P, Q, R.
    
        Has Related Work; no discussion section.
    
    * https://dl.acm.org/doi/epdf/10.1145/3719027.3765148
        Discussion, Related Work, and Conclusion.
        Conclusion is one paragraph: There's a gap. To address it we present X. We implemented it in/as Y. Our results show benefit Z.
    
    * https://dl.acm.org/doi/epdf/10.1145/3719027.3765116
        Has Limitations and Future Work section instead of Discussion. Has Related Works. Single paragraph conclusion.
        Conclusion: This paper proposed X, a new variant of Y that does Z. We provide A to do B and C. The verbed nouns P to achieve Q.
}

\comment{
    \pname{} offers the first binary CFI rewriter whose implementation is formally verified to produce policy-adhering binaries.
    Its proof of correctness relied on \Picinae{}, a bottom-up binary analysis platform in Rocq capable of reasoning about non-concrete binary code.
    The rewriter design was guided by several design policies for making it easier to formally verify, such policies include immutable security state, global state independence, and minimizing conditionals.
    The program-proof co-development cycle yielded several critical logic bugs that were revealed only because they halted the proof effort.
    That is, we claim no reasonable dynamic testing suite would have caught the bugs prior to deployment.
    
    This work also highlights the trusted components left to secure.
    % Among these is the internal lifter, used to translate code bytes to their computational semantics in PIL.
    This includes the internal lifter and the \ELFPackager.
    The internal lifter is ISA specific, and so far we only support ARMv7 and RISC-V.
    Creating a lifter for another ISA requires formalizing its instruction decoders and instruction semantics in Gallina, an especially daunting task for CISC architectures.
    The \ELFPackager combines the secured text section with the original ELF executable.
    At present, its correctness specification is impossible to formulate as it requires a formal correct implementation of a system loader, which has not been achieved.
    
    Future directions can also expand on the binary features supported by CFI (\eg system calls, dynamic linking) and apply this technique to other binary rewriting applications (\eg instrumentation, dynamic patching).
}

